\chapter{RELATED WORK}\label{chap_trabalhos_relacionados}

This chapter analyzes the state of the art in quantum machine learning, focusing on both theory and practical applications. The references discussed here were selected based on their relevance and contribution to the development of the proposed topic.

\section{Quantum Computing}

\citeonline{Mermin2007QComputerScience} provides a comprehensive introduction to quantum computing, discussing concepts such as qubits and cbits, the preparation and measurement of quantum states, and solutions to problems such as those of Deutsch, Bernstein--Vazirani, and Simon. The book also explores RSA cryptography and how to break it, search on quantum computers, and algorithms such as Grover iteration, in addition to addressing quantum error correction.

\citeonline{Wong2022IntroductionClassicalQuantum} presents a simplified introduction to classical and quantum computing. It covers basic concepts of classical computing, including bits, logic gates, and the programming of classical circuits using hardware description languages. The book also discusses classes of problems that are easy or difficult for computers and the prospect that quantum computers may be significantly faster than classical computers for some tasks. It then addresses the basic concepts of quantum computing, including qubits, superpositions, and quantum gates. The book emphasizes the importance of linear algebra in quantum computing and shows how it can be used. Ultimately, it aims to provide beginners with an introduction to the field of quantum computing.

\citeonline{Zubairy2020} presents the basic concepts of quantum mechanics to beginners and applies them to areas such as quantum communication and quantum computing. The book also discusses the foundations of quantum mechanics and their relationship with Newtonian dynamics. Its objective is to provide a comprehensive understanding of quantum mechanics and its applications.

\section{Quantum Machine Learning}

\citeonline{Biamonte2017QML} explores the intersection of quantum computing and machine learning and how they can be used together to improve learning algorithms and provide speedups. The article discusses several examples of how quantum mechanics can improve classical learning methods and how quantum algorithms can be used for various tasks. The authors also discuss the potential for quantum computing and machine learning to coevolve and become enabling technologies for one another.

\citeonline{Khan2020ML} provides an overview of quantum machine learning in comparison with classical approaches. The article discusses the technical contributions, strengths, and similarities of research in this domain. It also addresses recent progress in different quantum machine learning approaches, their complexity, and their applications in fields such as physics, chemistry, and natural language processing. The article explores the use of quantum systems to process classical data with machine learning algorithms and how machine learning algorithms and software can be formulated and implemented on quantum computers. It also highlights the significant hardware and software challenges that must be resolved before quantum machine learning becomes practical.

\citeonline{Havlicek2019qSupervisedLearning} presents a new method for supervised learning using enhanced quantum feature spaces. The authors propose a quantum-kernel-based approach that can classify datasets with high accuracy. They demonstrate the effectiveness of their method by applying it to several benchmark datasets and comparing the results with classical machine learning algorithms.

\citeonline{Rebentrost2014QSVM} demonstrates that the support vector machine, a popular classification method in machine learning, can be implemented on a quantum computer with logarithmic complexity in the feature-space dimension and the number of training data points. The article shows that this quantum algorithm can provide an exponential speedup over classical sampling algorithms, provided that the training data can be accessed in quantum superposition. It also highlights the data-privacy benefits of quantum machine learning and suggests that the quantum support vector machine could be used as a component of a larger quantum neural network.

\citeonline{Zhou2018HandwrittenDigitRecognition} develops a machine learning model for handwritten digit recognition based on a support vector machine and analyzes the influence of the sample size, kernel-function parameters, penalty coefficients, and other parameters on the prediction model. The study seeks to improve machine learning performance in the recognition of handwritten numerical information.

The book by \citeonline{Schuld2021MLQComputers} provides an introduction to and context for quantum machine learning. Its objective is to guide readers through different ways of thinking about machine learning from a quantum-computing perspective. The book also presents an example of how quantum computers can learn from data, illustrating several issues discussed throughout the text.

\citeonline{Métricas} discusses several evaluation metrics that data scientists can use to assess classification models and provides guidance on how and when to use them. The article emphasizes that different business problems require different optimization strategies and that accuracy is not always the best metric for evaluating model performance.

\section{Chapter Summary}

The related work described in this chapter was used to acquire theoretical knowledge of quantum physics, quantum computing, quantum machine learning methods, and their applications. Most of the references address quantum machine learning, while two books provide an introduction to quantum computing and clarify its physical foundations.